\question[5]
Die Klasse Parkhaus hat die Attribute plaetze (Anzahl aller Plätze) und belegt (Anzahl der belegten Plätze).
Dem Konstruktor wird die Anzahl der Plätze übergeben, zu Beginn ist kein Platz belegt.
\begin{itemize}
\item \texttt{frei()} gibt die Anzahl der freien Plätze zurück.
\item \texttt{einfahren()} belegt einen Platz und gibt True zurück. Ist kein Platz frei, wird False zurückgegeben.
\item \texttt{ausfahren()} gibt einen Platz wieder frei, wenn mindestens ein Platz belegt ist.
\end{itemize}
Implementiere die Klasse Parkhaus. Folgender Dialog soll möglich sein:
\begin{lstlisting}
p = Parkhaus(2)
print(p.einfahren(), p.einfahren(), p.einfahren())
print(p.frei())
p.ausfahren()
print(p)

True True False
0
1 von 2 Plätzen belegt
\end{lstlisting}
\begin{solutionbox}{9cm}
\begin{lstlisting}
class Parkhaus:
    def __init__(self, plaetze):
        self.plaetze = plaetze
        self.belegt = 0

    def frei(self):
        return self.plaetze - self.belegt

    def einfahren(self):
        if self.frei() > 0:
            self.belegt += 1
            return True
        return False

    def ausfahren(self):
        if self.belegt > 0:
            self.belegt -= 1

    def __str__(self):
        return f'{self.belegt} von {self.plaetze} Plätzen belegt'
\end{lstlisting}
\end{solutionbox}
